% Part: proof-theory
% Chapter: normalization
% Section: reductions

\documentclass[../../../include/open-logic-section]{subfiles}

\begin{document}

\olfileid{pt}{nor}{red}

\olsection{归约转换}

如果在 \Log{N1i} 的一个推导中，一个切割的长度为~$1$，那么它是一个由单个公式出现构成的片段，并且该公式同时是一个引入规则的结论和一个消去规则的主前提。在正规化证明中，这样的切割通过以下方式被“归约”：将那个以这样一个引入/消去配对结尾的子推导替换为一个新的推导，在新的推导中这对推理已被消除。这样的替换可能导致新的切割。然而，我们可以确认，新切割的秩低于被归约的那个切割。因此，新的推导要么具有更低的切割长度（因为一个长度为~$1$ 的极大切割已被移除），要么推导的切割秩已被降低（如果该极大切割是唯一的极大切割）。

例如，如果涉及的推理规则是 \Intro{\lif} 和
\Elim{\lif}，$!A \ident !B \lif !C$，并且存在一个以如下公式结尾的子推导~$\delta_1$：
\[
\AxiomC{$\Discharge{!B}{x}$}
\RightLabel{$\delta_2$}
\DeduceC{$!C$}
\DischargeRule{\Intro{\lif}}{x}
\UnaryInfC{$!B \lif !C$}
\AxiomC{}
\RightLabel{$\delta_3$}
\DeduceC{$!B$}
\RightLabel{\Elim{\lif}}
\BinaryInfC{$!C$}
\DisplayProof
\]
为了移除这个切割，
我们将~$\delta$ 中的子推导~$\delta_1$ 替换为以下内容：
\[
\AxiomC{}
\RightLabel{$\delta_3$}
\DeduceC{$!B$}
\RightLabel{$\Subst{\delta_2}{\delta_3}{!B^x}$}
\DeduceC{$!C$}
\DisplayProof
\]
$\Subst{\delta_2}{\delta_3}{!B^x}$ 指的是将~$\delta_2$ 中每一个带标签~$x$ 的假设~$!B$ 替换为整个推导~$\delta_3$ 所得到的推导。
这个推导不再有形如~$!B$ 且带有标签~$x$ 的未解除假设。
它也没有其他带标签~$x$ 的假设，因为与通常一样，我们假设没有两个公式作为假设出现时具有相同的标签。
因此，新的推导与~$\delta_1$ 具有相同的未解除假设，$\delta_2$ 中每一个解除假设的推理在~$\Subst{\delta_2}{\delta_3}{!B^x}$ 中也解除同样的假设，
而 $\delta_3$ 中每一个解除假设的推理，在~$\Subst{\delta_2}{\delta_3}{!B^x}$ 中则可能产生多个副本，这些副本解除了相应的假设。

任何完全位于 $\delta_2$、$\delta_3$ 或 $\delta_4$ 内的切割，或者完全位于~$\delta$ 中~$\delta_1$ 之外的切割都保持不变：
它们仍然涉及相同的公式，经历相同的推理步骤，因此特别地，其秩和长度与~$\delta$ 中相应的切割相同。
我们移除了一个长度为~$1$ 的极大切割，因此要么减少了切割长度，要么减少了切割秩（如果~$\delta$ 的切割长度是~$1$ 且这是唯一的极大切割）。

然而，可能发生以下两种情况：
\begin{enumerate}
  \item 通过将所有假设 $!B^x$ 替换为 $\delta_3$，我们倍增了~$\delta_3$ 中的所有切割。
  如果其中有些是极大切割，那么新推导的切割长度可能比旧推导的切割长度更大。
  我们必须保证这种情况不会发生。我们通过只对\emph{最右端}的极大切割应用归约转换来实现这一点，
  这里的“最右端”指的是：在穿过~$\delta$ 到达被归约切割的分支上，没有其他极大切割位于该切割的右边。
  如果~$\delta_3$ 中存在切割，它会在当前被归约切割的右边，那么我们就不是在处理一个最右端的切割。
  通过总是选择最右端的切割，我们可以保证：或者减少了推导的切割长度，或者甚至降低了切割秩。
  \item 如果 $\delta_2$ 中的一个假设 $!B^x$ 是某个消去规则的主前提，
  而~$\delta_3$ 的最后一个推理是 \Elim{\lor}、\Elim{\lexists} 或某个引入规则，
  那么通过将~$\delta_3$ 嫁接到 $\delta_2$ 的这样一个假设上，我们就生成了一个新的切割。
  原来在~$\delta_2$ 中仅仅是一个假设的地方，现在变成了一个长度为~$1$ 的切割（引入规则的结论且消去规则的主前提），
  或者变成了一个切割段的最后一个公式出现（消去规则的主前提且 \Elim{\lor} 或 \Elim{\lexists} 的结论）。
  如果 $B^x$ 是 \Elim{\lor} 或 \Elim{\lexists} 的次前提，它原本已经是一个段的首个公式出现，并且可能属于一个切割段。
  在新的推导中，这个段现在可能变长了。
  然而，注意构成这个新切割或现有切割段的公式是~$!B$，并且
  $\depth{!B} < \depth{!B \lif !C}$。因此，尽管我们可能增加了切割或增大了切割的长度，
  但这些切割都不是极大切割。所以，我们要么减少了切割秩，
  要么，如果有相同秩的切割仍然存在，则减少了切割长度。
\end{enumerate}

这种减少长度为~$1$ 的切割的过程称为\emph{归约转换}（或称\emph{绕路转换}）。
部分规则对对应的转换模式见 \olref{tab:reduction-conversions}。（其余的留作练习。） 

\begin{table}
\begin{multline*}
\bottomAlignProof
\AxiomC{$\Discharge{!B}{x}$}
\RightLabel{$\delta_2$}
\shortDeduce
\DeduceC{$!C$}
\DischargeRule{\Intro{\lif}}{x}
\UnaryInfC{$!B \lif !C$}
\AxiomC{}
\RightLabel{$\delta_3$}
\shortDeduce
\DeduceC{$!B$}
\RightLabel{\Elim{\lif}}
\BinaryInfC{$!C$}
\DisplayProof
\quad \longrightarrow\quad
\shoveright{\bottomAlignProof
\AxiomC{}
\RightLabel{$\delta_3$}
\shortDeduce
\DeduceC{$!B$}
\RightLabel{$\Subst{\delta_2}{\delta_3}{!B^x}$}
\shortDeduce
\DeduceC{$!C$}
\DisplayProof}
\\[4ex]
\shoveleft{\bottomAlignProof
\AxiomC{}
\RightLabel{$\delta_2$}
\shortDeduce
\DeduceC{$!B$}
\AxiomC{}
\RightLabel{$\delta_3$}
\shortDeduce
\DeduceC{$!C$}
\RightLabel{\Intro{\land}}
\BinaryInfC{$!B \land !C$}
\RightLabel{\Elim{\land}[1]}
\UnaryInfC{$!B$}
\DisplayProof
\quad \longrightarrow\quad
\bottomAlignProof
\AxiomC{}
\RightLabel{$\delta_2$}
\shortDeduce
\DeduceC{$!B$}
\DisplayProof}\qquad
\shoveright{\bottomAlignProof
\AxiomC{}
\RightLabel{$\delta_2$}
\shortDeduce
\DeduceC{$!B$}
\AxiomC{}
\RightLabel{$\delta_3$}
\shortDeduce
\DeduceC{$!C$}
\RightLabel{\Intro{\land}}
\BinaryInfC{$!B \land !C$}
\RightLabel{\Elim{\land}[2]}
\UnaryInfC{$!C$}
\DisplayProof
\quad \longrightarrow\quad
\bottomAlignProof
\AxiomC{}
\RightLabel{$\delta_3$}
\shortDeduce
\DeduceC{$!C$}
\DisplayProof}
\\[4ex]
\bottomAlignProof
\AxiomC{}
\RightLabel{$\delta_2$}
\shortDeduce
\DeduceC{$!B$}
\RightLabel{\Intro{\lor}[1]}
\UnaryInfC{$!B \lor !C$}
\AxiomC{$\Discharge{!B}{x}$}
\RightLabel{$\delta_3$}
\shortDeduce
\DeduceC{$!D$}
\AxiomC{$\Discharge{!C}{y}$}
\RightLabel{$\delta_4$}
\shortDeduce
\DeduceC{$!D$}
\DischargeRule{\Elim{\lor}}{x\,y}
\TrinaryInfC{$!D$}
\DisplayProof
\quad \longrightarrow\quad
\shoveright{\bottomAlignProof
\AxiomC{}
\RightLabel{$\delta_2$}
\shortDeduce
\DeduceC{$!B$}
\RightLabel{$\Subst{\delta_3}{\delta_2}{!B^x}$}
\shortDeduce
\DeduceC{$!D$}
\DisplayProof}
\\[2ex]
\shoveleft{\bottomAlignProof
\AxiomC{}
\RightLabel{$\delta_2$}
\shortDeduce
\DeduceC{$!B(c)$}
\RightLabel{\Intro{\lforall}}
\UnaryInfC{$\lforall[x][!B(x)]$}
\RightLabel{\Elim{\lforall}}
\UnaryInfC{$!B(t)$}
\DisplayProof
\quad \longrightarrow\quad
\bottomAlignProof
\AxiomC{}
\RightLabel{$\Subst{\delta_2}{t}{c}$}
\shortDeduce
\DeduceC{$!B(t)$}
\DisplayProof}
\\[4ex]
\shoveright{\bottomAlignProof
\AxiomC{}
\RightLabel{$\delta_2$}
\shortDeduce
\DeduceC{$!B(t)$}
\RightLabel{\Intro{\lexists}}
\UnaryInfC{$\lexists[x][!B(x)]$}
\AxiomC{$\Discharge{!B(c)}{x}$}
\RightLabel{$\delta_3$}
\shortDeduce
\DeduceC{$!D$}
\DischargeRule{\Elim{\lexists}}{x}
\BinaryInfC{$!D$}
\DisplayProof
\quad \longrightarrow\quad
\bottomAlignProof
\AxiomC{}
\RightLabel{$\delta_2$}
\shortDeduce
\DeduceC{$!B(t)$}
\RightLabel{$\Subst{\Subst{\delta_3}{t}{c}}{\delta_2}{!B(t)^x}$}
\shortDeduce
\DeduceC{$!D$}
\DisplayProof}
\end{multline*}
\caption{\Log{N1i} 的一些归约转换。}
\ollabel{tab:reduction-conversions}
\end{table}

\begin{prob}
给出针对 \Intro{\lnot} 紧接 \Elim{\lnot} 规则的归约转换。
\end{prob}

还有一种情况需要考虑。切割的定义包含了切割长度为~$1$ 且 $!A$ 既是~\FalseInt{} 的结论又是某个消去规则的主前提的情形。
这种情况的处理方式类似于 $!A$ 是某引入规则结论时的情形。
例如，将
\[
\bottomAlignProof
\AxiomC{}
\RightLabel{$\delta_2$}
\DeduceC{$\lfalse$}
\UnaryInfC{$!B \lif !C$}
\AxiomC{}
\RightLabel{$\delta_3$}
\DeduceC{$!B$}
\RightLabel{\Elim{\lif}}
\BinaryInfC{$!C$}
\DisplayProof
\quad 替换为 \quad
\bottomAlignProof
\AxiomC{}
\RightLabel{$\delta_2$}
\DeduceC{$\lfalse$}
\UnaryInfC{$!C$}
\DisplayProof
\]
对于 $!A \ident (!B \lor !C)$ 或 $!A \ident \lexists[x][!B(x)]$ 的情况，我们必须稍加小心。
例如，如果我们把
\[
\bottomAlignProof
\AxiomC{}
\RightLabel{$\delta_2$}
\DeduceC{$\lfalse$}
\RightLabel{\FalseInt}
\UnaryInfC{$!B \lor !C$}
\AxiomC{$\Discharge{!B}{x}$}
\RightLabel{$\delta_3$}
\DeduceC{$!D$}
\AxiomC{$\Discharge{!C}{x}$}
\RightLabel{$\delta_4$}
\DeduceC{$!D$}
\RightLabel{\Elim{\lor}}
\TrinaryInfC{$!D$}
\DisplayProof
\quad 替换为 \quad
\bottomAlignProof
\AxiomC{}
\RightLabel{$\delta_2$}
\DeduceC{$\lfalse$}
\RightLabel{\FalseInt}
\UnaryInfC{$!D$}
\DisplayProof
\]
就可能引入一个以~$!D$ 为切割公式的新切割，并且无法保证 $\depth{!D} < \depth{!B \lor !C}$!{}
所以，我们必须像处理 \Intro{\lor}/\Elim{\lor} 归约那样进行替换，即用
\[
\bottomAlignProof
\AxiomC{}
\RightLabel{$\delta_2$}
\DeduceC{$\lfalse$}
\RightLabel{\FalseInt}
\UnaryInfC{$!B$}
\RightLabel{$\delta_3$}
\DeduceC{$!D$}
\DisplayProof
\quad 或 \quad
\bottomAlignProof
\AxiomC{}
\RightLabel{$\delta_2$}
\DeduceC{$\lfalse$}
\RightLabel{\FalseInt}
\UnaryInfC{$!C$}
\RightLabel{$\delta_4$}
\DeduceC{$!D$}
\DisplayProof
\] 来替换

\end{document}